% Copy-ready replacement for src/core.tex, Proposition Functorial modulus maps.
% This private correction does not alter the supplied archive.
\begin{proposition}[Functorial modulus maps and their compactifications]
\label{prop:crt}
For positive odd integers $D\mid E$, coordinate reduction
$\mathcal F_E\to\mathcal F_D$ is equivariant and surjective, with fibres
of size $(E/D)^3$. It induces a map of the corresponding compactified
covers. At a marked source branch cycle of length $e_E$ mapping to a
target branch cycle of length $e_D$, one has $e_D\mid e_E$, and the local
degree is $e_E/e_D$. At the cusp this is $E/D$.

For coprime positive odd integers $D,E$, coordinate CRT gives an
equivariant bijection
\[
 \mathcal F_{DE}\xrightarrow{\sim}\mathcal F_D\times\mathcal F_E,
\]
with coordinate inverse
$a+D((b-a)D^{-1}\bmod E)\bmod DE$.
The corresponding unramified cover is the fibre product over the
punctured base. Its compactification is the normalization of the fibre
product of the compactified covers over $\mathbb P^1$.
The cocycles and local systems have the induced reduction and CRT maps.
Every product and every map retains the same base and candidate label.
\end{proposition}
\begin{proof}
All three affine formulas have integral coefficients, hence commute with
coordinate reduction. Each residue modulo $D$ has exactly $E/D$ lifts
modulo $E$, proving the fibre count in three coordinates. Equivariance
gives the map of covers away from the punctures.

At a fixed puncture, let $e_E$ be the source cycle length and $e_D$ its
image cycle length. The $e_E$-th power of target monodromy fixes the
image sheet, so $e_D\mid e_E$. Choose compatible sheet origins and
local coordinates with base parameter
\[
 t=w_E^{e_E}=w_D^{e_D}.
\]
The prescribed punctured-cover map has the local form
\[
 w_D=\zeta w_E^{e_E/e_D},\qquad \zeta^{e_D}=1,
\]
where the root of unity retains the target branch mark. The positive
integer exponent proves holomorphic extension, unique by agreement on
the punctured disc. At infinity $e_E=E$ and $e_D=D$ by the cusp-cycle
theorem. At the finite branch points their values are the actual cycle
lengths dividing three and four; they need not equal the moduli.

The displayed CRT formula is inverse to both coordinate reductions.
It commutes with all three monodromies, so it identifies the unramified
level-$DE$ cover with the fibre product of the two unramified covers.
To complete this identification, consider source and target branch
cycles of lengths $e_D,e_E$. Their local fibre product has equation
\[
 x^{e_D}=y^{e_E}.
\]
Put $d=\gcd(e_D,e_E)$, $r=e_D/d$, $s=e_E/d$. Its $d$ branches are
$x^r=\xi y^s$ for $\xi^d=1$. On a branch, choose $\eta^r=\xi$.
The normalization has the parametrization
\[
 x=\eta w^s,\qquad y=w^r,\qquad
 t=w^{\operatorname{lcm}(e_D,e_E)}.
\]
The two maps to the base agree because $\eta^{e_D}=\xi^d=1$.
Since $\gcd(r,s)=1$, Bezout's identity gives a birational inverse in
the field of meromorphic germs, and the parameter disc is smooth.
The product monodromy has precisely $d$ cycles of length
$\operatorname{lcm}(e_D,e_E)$: a simultaneous iterate returns exactly
when divisible by both lengths, and the number of orbits is
$e_De_E/\operatorname{lcm}(e_D,e_E)=d$.
Thus the normalized local branches are exactly the completed product
cover. Their unique extensions glue to the compactified level-$DE$
cover. This argument is componentwise for disconnected covers and
keeps every candidate label and marked branch.

For example, reduction from level 15 to level 3 sends the $a_1$ cycle
\[
 (0,0,0)\mapsto(6,9,13)\mapsto(0,9,2)\mapsto(0,0,0)
\]
to the cycle
\[
 (0,0,0)\mapsto(0,0,1)\mapsto(0,0,2)\mapsto(0,0,0).
\]
The local degree here is $3/3=1$, whereas at the cusp it is $15/3=5$.
For the coprime levels three and five the cusp fibre product
$x^3=y^5$ is singular; its normalization is
$x=w^5,y=w^3,t=w^{15}$. This records why normalization is part of
the compactified product statement.
\end{proof}
